% TOP_PROBLEM: {"id": 2307056, "title": "Research Problems in Function Theory — Problem 7.56", "category_id": 8, "category": {"id": 8, "name": "analysis", "display_name": "Analysis", "description": "Limits, continuity, calculus, and function theory.", "slug": "analysis", "order_index": 8, "created_at": "2026-07-31T15:26:25.671Z"}, "status": "open"}
% FIRST_POSTED: null
% ENTRY_KIND: statement_only
% Imported from: batch06.tex; UnsolvedMath ID 2307056
% Compile twice: lualatex 2307056.tex
\documentclass[11pt]{article}
\usepackage{fontspec}
\setmainfont{Latin Modern Roman}
\usepackage{amsmath,amssymb,mathtools,unicode-math}
\setmathfont{Latin Modern Math}
\usepackage{xurl,hyperref}
\hypersetup{hidelinks,pdftitle={UnsolvedMath Problem 2307056}}
\newcommand{\sref}[2]{\hyperlink{source:#1:#2}{[S#2]}}
\newcommand{\cataloguescope}[1]{\paragraph{Mathematical question.}#1}
\begin{document}
% UnsolvedMath ID 2307056; source catalogue code AMR-022-7056
\setcounter{section}{2307055}
\section{Research Problems in Function Theory — Problem 7.56}\label{2307056}
{\small\sffamily UnsolvedMath ID 2307056 · Analysis}
\subsection{Definitions and mathematical statement}

\paragraph{Shared notation.}$\mathbb N=\{1,2,\ldots\}$, and $\mathbb Z,\mathbb Q,\mathbb R,\mathbb C$ denote the integers, rationals, reals and complex numbers. Any other conventions are specified locally.


Let $\widehat{\mathbb C}=\mathbb C\cup\{\infty\}$, and let $\Gamma\subset\widehat{\mathbb C}$ be a Jordan curve through $\infty$. Write $H_+=\{\operatorname{Im}z>0\}$ and $H_-=\{\operatorname{Im}z<0\}$. Let conformal bijections $f_\pm:H_\pm\to\Omega_\pm$ map onto the two complementary domains of $\Gamma$, with boundary extensions satisfying $f_\pm(\infty)=\infty$. Their boundary identification is
\[h=(f_-|_{\mathbb R})^{-1}\circ(f_+|_{\mathbb R}):\mathbb R\to\mathbb R.
\]
It is an increasing homeomorphism. Quasisymmetry means that some $C\ge1$ satisfies $C^{-1}\le(h(x+t)-h(x))/(h(x)-h(x-t))\le C$ for all $x\in\mathbb R$ and $t>0$.
\paragraph{Statement recorded in the supplied source.}
Characterize the homeomorphisms $h$ arising by these conformal boundary identifications for arbitrary Jordan curves. In particular, does every increasing homeomorphism of $\mathbb R$ arise? Quasisymmetric identifications correspond to quasicircles. The original update records Oikawa and Huber counterexamples to the every-homeomorphism assertion. \sref{2307056}{2}
\subsection{Short English statement}
Characterize the boundary homeomorphisms obtained by conformally welding the two sides of a Jordan curve, including whether every increasing homeomorphism occurs.
\subsection{Sources}
\begin{description}
\item[\hypertarget{source:2307056:1}{S1}] ULAM.AI and the UnsolvedMath Contributors, UnsolvedMath: A Curated Collection of Open Mathematics Problems, record 2307056 (AMR-022-7056). CC BY 4.0. \url{https://huggingface.co/datasets/ulamai/UnsolvedMath}.
\item[\hypertarget{source:2307056:2}{S2}] Walter K. Hayman and Edward F. Lingham, Research Problems in Function Theory (2018), Problem 7.56, its complete update and relevant preceding definitions. \url{https://arxiv.org/abs/1809.07200}.
\end{description}
\label{end:2307056}
\end{document}
